\section{Motivation}

\subsection{Repository Context Repeats Across Agents}

In a coding MAS workflow, agents have different roles but often consult the same code base.
For example, a Planner may receive a repository summary and several large files to decide a repair strategy.
An Implementer later receives the issue, the plan, and two of the same files.
A Debugger may receive test output plus one of the same files.
The repeated code is exact, but its prompt position changes because each agent has a different prefix and context.
This is the common case rather than an artifact of a particular prompt template: role instructions, chain-of-thought summaries, test output, and tool results are usually prepended before the code that a later agent must inspect.
As a result, the expensive tokens are stable as repository content but unstable as prefix positions.

This pattern is not rare in repo-level tasks.
Our 500-case SWE-bench Verified code-base-content manifest contains 1500 large Python file segments, 5.97M source lines, and an estimated 58.8M reusable tokens.
The repeated data is large enough to matter for serving latency and memory movement, yet structured enough for templates to describe future accesses.
The serving problem is therefore not simply ``can we cache a prompt?'' but ``can we preserve useful KV for code objects that recur across different agents?''
Treating repository files as first-class objects gives the runtime a stable name and signature for content that would otherwise appear only as a long substring inside each request.

\subsection{Agent Templates Expose Serving Semantics}

The previous project draft, AgentTemplateKV, observed that multi-agent code programs are not arbitrary conversations.
They have repeatable task-specific DAGs such as plan--implement--test--debug--review, and edges carry typed content such as patches, code files, test logs, or review comments.
\system keeps this template view but narrows the serving objective.
Rather than assigning generic TTLs to every agent turn, \system asks which exact code-base segments will recur in later agents and when the engine should prefetch or retain their KV.

The template side therefore records role, downstream consumers, edge content type, and use distance.
These fields are compiled into \prefetchhint and priority metadata.
The engine can use them for prefetch and eviction without trusting them for correctness.
Correctness is still determined by the exact content gate at reuse time.
This separation is important because templates are predictive.
They may be incomplete, stale, or conservative; a later agent may skip a file, add a new diagnostic log, or receive a shortened prompt after a tool failure.
\system treats such deviations as scheduling misses rather than correctness failures.
A wrong hint can waste prefetch work, but it cannot authorize reuse of non-identical code.

\subsection{Prefix Caches Are Not Enough}

Prefix caching works when requests share the same initial token sequence.
Coding MAS prompts can share a system prefix, but the expensive repository content is often placed after role-specific instructions or intermediate outputs.
Two prompts may therefore contain the same file at positions $i$ and $j$.
The cached K/V tensors for the span are not position-free: RoPE encodes positions into attention states.
This makes direct middle-span reuse numerically different from a fresh prefill.
One conservative solution is to recompute every non-prefix code token for every agent.
That solution is simple but leaves substantial reuse opportunity unused, especially when two or three agents inspect the same thousand-line file.
The alternative explored here is narrower: reuse only when the code content is byte-for-byte equivalent after normalization, and then repair the positional mismatch with a RoPE delta transform.
This keeps the semantic risk focused on numerical approximation rather than on approximate code matching.

\subsection{Syntax Is a Locator, Not a Safety Gate}

AST and anchor metadata help find corresponding code regions.
They are not sufficient to prove safe reuse.
Two functions can share names or syntax shape while changing a literal, operator, or call target.
Our controlled gate suite confirms this: AST-only, span-overlap-only, and no-gate policies falsely accept changed code, while exact-content gating rejects every non-identical case.
This observation drives the central safety invariant of \system: only exact code content can be reused.
The point is not that syntax metadata is useless.
On the contrary, it makes candidate lookup practical by narrowing the search to likely spans and by attaching source-level explanations to cache entries.
The safety boundary is simply placed after lookup.
Locators answer where a reusable object might be; signatures answer whether the cached object is the same code.
